\documentclass[a4paper,11pt]{article}
\usepackage[utf8]{inputenc}
\usepackage[T1]{fontenc}
\usepackage[english]{babel}
\usepackage[left=2.5cm,right=2.5cm,top=2cm,bottom=2cm]{geometry}
\usepackage{hyperref}
\usepackage{xcolor}
\usepackage{microtype}
\usepackage{lmodern}

\definecolor{primary}{RGB}{0, 51, 102}
\hypersetup{colorlinks=true, urlcolor=primary}

\begin{document}
\noindent \textbf{Lo\"ic Aron MBASSI EWOLO} \\
ENSEA -- Double Dipl\^ome Ing\'enieur \\
\href{mailto:loicaron.mbassiewolo@ensea.fr}{loicaron.mbassiewolo@ensea.fr} \quad (+33) 7 59 52 33 98 \\
\href{https://github.com/Nameless0l}{github.com/Nameless0l} \quad \href{https://mbassiloic.tech}{mbassiloic.tech}

\vspace{1em}\noindent Saint-Germain-en-Laye, le \today
\vspace{1em}\noindent \textbf{Objet : Candidature Alternance -- Ing\'enieur D\'eveloppement Logiciel Embarqu\'e C/C++}
\vspace{1.5em}\noindent Madame, Monsieur,

\vspace{0.8em}\noindent J'ai d\'evelopp\'e des logiciels embarqu\'es complets en C/C++ sur STM32 avec interfaçage de capteurs IMU et Flex via protocoles I2C/SPI dans le projet Harmony Gloves (Labo IOTA ENSEA), et impl\'ement\'e un syst\`eme SoPC VHDL avec soft-processeur Nios V, debug JTAG et contr\^oleur I2C en cours \`a l'ENSEA. Formation double dipl\^ome ENSEA/Polytechnique Yaound\'e, sp\'ecialit\'e syst\`emes embarqu\'es et \'electronique temps r\'eel. Stage INRIA Saclay en cours sur l'analyse statique de syst\`emes logiciels embarqu\'es.

\vspace{0.8em}\noindent Exail, sp\'ecialiste mondial de la navigation inertielle et de la robotique autonome, repr\'esente un environnement id\'eal pour d\'evelopper des logiciels temps r\'eel critiques sur architectures ARM pour centrales inertielles et syst\`emes de guidage. Mon exp\'erience de protocoles bas niveau (I2C, SPI, UART) et validation oscilloscope m'a permis de maîtriser les contraintes temps r\'eel strictes requises en navigation inertielle.

\vspace{0.8em}\noindent Mon exp\'erience Jetson CoHoMa (Challenge D\'efense ENSEA/Arm\'ee) valide ma capacit\'e \`a d\'evelopper des logiciels C/C++ temps r\'eel optimis\'es sur architectures ARM pour robotique autonome, avec gestion m\'emoire limit\'ee et parallélisation. J'ai aussi implement\'e des syst\`emes multiprocessus robustes (PROJET-TAXI) avec IPC, synchronisation et gestion de signaux Linux, une exp\'erience transposable aux d\'efis d'OS temps r\'eel comme QNX.

\vspace{0.8em}\noindent Je suis disponible en alternance d\`es septembre 2026 pour une dur\'ee de 12 mois, mobile en Île-de-France. Je serais ravi de discuter de comment mes comp\'etences en logiciel embarqu\'e C/C++, temps r\'eel critique et interfaçage hardware peuvent contribuer aux syst\`emes de navigation inertielle d'Exail pour les domaines maritime, a\'erien et robotique.

\vspace{1.5em}\noindent Veuillez agr\'eer, Madame, Monsieur, l'expression de mes salutations distingu\'ees.
\vspace{2em}\noindent \textbf{Lo\"ic Aron MBASSI EWOLO}
\end{document}
